\abstract{

We establish new bounds on the Grothendieck constant \(\KG\):
\[
    \frac{6\pi}{11}
    \le
    \KG
    \le
    \frac{\pi}{2\log(1+\sqrt2)} - 10^{-4}.
\]
Methodologically, our lower bound approach differs from previous works by establishing limitations on the asymptotically optimal Krivine schemes, rather than giving explicit constructions of gap instances.
Our upper bound is obtained by proposing and analyzing the first asymptotic construction of rounding schemes, whereas previous works only consider low-dimensional schemes.
Together, these bounds determine the previously unknown tenths digit of \(\KG\) to be $7$. 
The bounds were discovered by a long-running collaborative effort of humans and a long-horizon AI research system that we engineered. A detailed account of the AI-assisted research process appears in a companion paper~\cite{lisaha2026longhorizon}.

% ~\footnote{\cite{lisaha2026longhorizon}}



% Both bounds were discovered through substantial human--AI collaboration; all stated theorems were independently verified by the authors.


% Methodologically, our lower bound approach differs from previous works by establishing limitations on Krivine schemes which are known to be asymptotically optimal, rather than giving an explicit construction of gap instances.
% The upper bound is obtained by proposing and analyzing the first asymptotic construction of rounding schemes.


% The lower bound is proved by establishing limitations within the Krivine rounding framework, while the upper bound is obtained from a new asymptotic family of rounding schemes.
% Together, these bounds determine the previously unknown tenths digit of \(\KG\) to be \(7\).

% Both bounds were discovered through substantial human--AI collaboration; all stated theorems were independently verified by the authors.


%%%%%%%%%%%%%%%%%%%%%%%%%%%%%%%%%%%%%%%%

% We improve the lower bound on the Grothendieck constant to
% \[
% \KG\ge \frac{6\pi}{11} = 1.7135 \dots
% \]
% Methodologically, our approach differs from previous works by establishing limitations on Krivine schemes which are known to be asymptotically optimal, rather than giving an explicit construction of gap instances.

% We also propose and analyze the first asymptotic construction of rounding schemes for the Grothendieck constant (\(\KG\)), improving the best upper bound,
% \[
%   \KG\le {\pi\over 2\log(1+\sqrt2)} - 3 .47 \times 10^{-4} = 1.7818 \ldots
% \]
% In particular, combining the two bounds shows the tenths digit of $\KG$ to be $7$.

% % {\bf \color{red} RS: Attribute human collaboration briefly}. Both the lower and upper bounds were found with substantial AI assistance, including a long-horizon AI research system.

% Both the lower and upper bounds were found through substantial human--AI collaboration, including the use of a long-horizon AI research system. All results stated as theorems were independently verified by the authors.


%%%%%%%%%%%%%%%%%%%%%%%%%%%%%%%%%%%%%%%%


% We analyze a family of high-dimensional Krivine rounding schemes.  
% The simplest member of the family uses an anti-aligned third-chaos average and an aligned fifth-chaos average.  The finite-dimensional schemes are ordinary Gaussian rounding schemes, while their limiting correlation retains the Hermite degrees through the covariance powers $t^3$ and $t^5$.
% We derive explicit coefficient formulas for the limiting normalized correlation $H(t)=\sum b_m t^m$, prove that strict limiting inverse-majorant certificates transfer to sufficiently high finite dimensions, and give an interval-arithmetic certificate showing that the nonlinear coefficient mass of $H$ is small.  The inverse-majorant step uses a Wiener-algebra contraction, and the coefficient tail is controlled by a scalar Sobolev norm $B_3=\norm{\Dop^3H}_{L^2(\T)}$ that is itself bounded by a certified two-dimensional Gaussian quadrature.
% For the parameters
% \[
  % \eta=0.136419125,\qquad s_3=0.34101124,\qquad s_5=0.05276111,
% \]
% Our construction gives a finite-dimensional Krivine rounding scheme with
% \[
  % \KG\le {\pi\over 2\gamma}=1.7818666069360661,
  % \qquad \gamma=0.881545409.
% \]
% Equivalently,
% \[
%   \KG\le {\pi\over 2\log(1+\sqrt2)}-\cdot10^{-3}.
% \]
}
\newpage 

\tableofcontents 

% \newpage